We model the genetic programming system as a plant and the outer Bayesian
optimization loop as a controller. At control step $\ell$, the controller
selects
%
\begin{equation}
u_{\ell} = (p_{c,\ell}, p_{m,\ell}, k_{\ell}),
\end{equation}
%
where $p_{c,\ell}$ is the crossover rate, $p_{m,\ell}$ is the mutation
rate, and $k_{\ell}$ is the number of generations until the next control
update.

If $g$ denotes the GP generation index, the next update point is given by
%
\begin{equation}
g_{\ell+1} = g_{\ell} + k_{\ell}.
\end{equation}
%
Let $P_{g}$ be the population at generation $g$, $F$ the GP transition,
and $\xi_{\ell}$ the stochastic disturbance. The population transition over
one control interval is written as
%
\begin{equation}
P_{g_{\ell+1}} = F(P_{g_{\ell}}, u_{\ell}, \xi_{\ell}).
\end{equation}

The controller observes a compact state representation,
%
\begin{equation}
x_{\ell} =
\left[
\tau_{\ell},
\widetilde{HV}_{\ell},
\widetilde{\Delta HV}_{\ell},
D_{\ell},
\widetilde{L}_{\ell},
\widetilde{s}_{\ell}
\right]^{\top},
\end{equation}
%
which summarizes the evolutionary stage, current performance, recent
improvement, diversity, average tree size, and stagnation length.

This section can be extended with the acquisition function, action-space
constraints, reward definition, and the update policy used by the
controller.
